\documentclass[11pt,a4paper]{article}
\usepackage[utf8]{inputenc}
\usepackage[T1]{fontenc}
\usepackage[catalan]{babel}
\usepackage{lmodern}
\usepackage[sfdefault,lining]{FiraSans}
\usepackage[a4paper,top=2.3cm,bottom=2.2cm,left=1.8cm,right=1.8cm,headheight=1.3cm,headsep=0.4cm,footskip=1cm]{geometry}
\usepackage{amsmath,amssymb}
\usepackage{icomma}
\usepackage{graphicx}
\usepackage[table]{xcolor}
\usepackage{tikz}
\usepackage{fontawesome5}
\usepackage{ccicons}
\usepackage{fancyhdr}
\usepackage{lastpage}
\usepackage{tasks}
\usepackage{needspace}
\usepackage[most]{tcolorbox}
\usepackage{array}
\usepackage[hidelinks,pdftitle={Fitxa d'exercicis de la unitat 3: Introducció a l'àlgebra. Polinomis. 2n d'ESO},pdfsubject={Matemàtiques 2n ESO},pdfkeywords={llenguatge algebraic, valor numèric, monomis, monomis semblants, operacions amb monomis, potència d'un monomi, polinomis, suma i resta de polinomis, producte de polinomis, identitats notables, factor comú}]{hyperref}

\definecolor{navy}{HTML}{1B2748}
\definecolor{accent}{HTML}{E8603A}
\definecolor{teal}{HTML}{1F9E90}
\definecolor{soft}{HTML}{EEF1F8}
\definecolor{muted}{HTML}{5D6685}

\setlength{\parindent}{0pt}
\setlength{\parskip}{3pt}
\settasks{label=\alph*), label-format=\bfseries\color{navy}, label-width=1.3em, item-indent=1.8em, column-sep=0.6em, after-item-skip=5pt, before-skip=3pt, after-skip=4pt}

% ---------- capçalera i peu ----------
\newcommand{\llicencia}{\href{https://creativecommons.org/licenses/by-nc-sa/4.0/deed.ca}{\ccbyncsa}\ \ Llicència Creative Commons Reconeixement-NoComercial-CompartirIgual 4.0 Internacional (CC BY-NC-SA 4.0)}
\pagestyle{fancy}
\fancyhf{}
\fancyhead[L]{\raisebox{-0.35\height}{\includegraphics[height=0.95cm]{logo-gva.png}}}
\fancyhead[C]{\small\color{muted}Fitxa d'exercicis · Unitat 3: Introducció a l'àlgebra. Polinomis · 2n d'ESO}
\fancyhead[R]{\raisebox{-0.35\height}{\includegraphics[height=0.95cm]{logo-vermellar.png}}}
\fancyfoot[L]{\scriptsize\color{muted}\llicencia}
\fancyfoot[R]{\small\color{muted}\thepage}
\renewcommand{\headrulewidth}{0.4pt}
\renewcommand{\headrule}{\color{muted!40}\hrule height \headrulewidth}
\fancypagestyle{primera}{\fancyhf{}\renewcommand{\headrulewidth}{0pt}
  \fancyfoot[L]{\scriptsize\color{muted}\llicencia}
  \fancyfoot[R]{\small\color{muted}\thepage}}

% ---------- nivells i exercicis ----------
\newcommand{\estrelles}[1]{{\color{accent}\ifcase#1\or\faStar\faStar[regular]\faStar[regular]\or\faStar\faStar\faStar[regular]\or\faStar\faStar\faStar\fi}}
\newcounter{exer}
\newcommand{\ex}[1]{\par\needspace{4\baselineskip}\medskip\stepcounter{exer}\noindent{\color{navy}\bfseries\theexer.}\ \ {\small\estrelles{#1}}\ \ }

\newcommand{\seccio}[2]{\par\needspace{10\baselineskip}\bigskip
  \begin{tcolorbox}[colback=navy,colframe=navy,coltext=white,arc=2mm,boxrule=0pt,left=3mm,right=3mm,top=1.2mm,bottom=1.2mm]
  {\large\bfseries #1\quad #2}
  \end{tcolorbox}\vspace{-2pt}}

\newcolumntype{C}[1]{>{\centering\arraybackslash}p{#1}}
% quadre de recordatori (només a la fitxa d'exercicis)
\newcommand{\recorda}[1]{\par\smallskip\begin{tcolorbox}[colback=accent!7,colframe=accent!55,arc=2mm,boxrule=0.5pt,left=2.5mm,right=2.5mm,top=1mm,bottom=1mm]\small{\color{accent}\faLightbulb[regular]}\ \textbf{Recorda:}\ #1\end{tcolorbox}}
\newcommand{\res}[1]{#1}\newcommand{\rest}[1]{#1}

\def\fmneed{12}
% =====================================================================
% Capçalera FentMates de les fitxes: un color per a cada curs.
% Els colors són els de la pàgina inicial (LEVELS a assets/js/materials.js);
% el logo ve de _eines/logo/pdf/ (python3 _eines/logo/genera.py).
%
% Ús, al final del preàmbul de la unitat:
%   % =====================================================================
% Capçalera FentMates de les fitxes: un color per a cada curs.
% Els colors són els de la pàgina inicial (LEVELS a assets/js/materials.js);
% el logo ve de _eines/logo/pdf/ (python3 _eines/logo/genera.py).
%
% Ús, al final del preàmbul de la unitat:
%   % =====================================================================
% Capçalera FentMates de les fitxes: un color per a cada curs.
% Els colors són els de la pàgina inicial (LEVELS a assets/js/materials.js);
% el logo ve de _eines/logo/pdf/ (python3 _eines/logo/genera.py).
%
% Ús, al final del preàmbul de la unitat:
%   \input{../comu/fentmates.tex}
%   \fmcurs{1eso}
% i, al començament del document:
%   \capcalera{Matemàtiques}{1r ESO}{Fitxa d'exercicis · Unitat 2}{Divisibilitat}
% Opcional abans de \input: \def\fmneed{14} (espai mínim abans de cada secció).
% =====================================================================
\makeatletter
\def\fm@logodir{../../logo/pdf}
\providecommand{\fmneed}{10}
\def\fm@cursnom{}

% clau del curs, color i color del text damunt del color
\newcommand{\fm@curs}[3]{\@namedef{fm@color@#1}{#2}\@namedef{fm@text@#1}{#3}}
\fm@curs{1eso}{E6B84A}{navy}
\fm@curs{2eso}{3B82F6}{white}
\fm@curs{3eso}{FF7A50}{white}
\fm@curs{4eso-a}{A78BFA}{white}
\fm@curs{4eso-b}{6D28D9}{white}
\fm@curs{1batx-mi}{9B1C1C}{white}
\fm@curs{1batx-ccss}{EF6B6B}{white}
\fm@curs{2batx-mii}{1E6B47}{white}
\fm@curs{2batx-ccss}{2BB3A3}{white}

\newcommand{\fmcurs}[1]{%
  \@ifundefined{fm@color@#1}{\PackageError{fentmates}{Curs desconegut: #1}{}}{}%
  \edef\fm@tmp{\noexpand\definecolor{curs}{HTML}{\@nameuse{fm@color@#1}}%
               \noexpand\colorlet{curstext}{\@nameuse{fm@text@#1}}}\fm@tmp
  \colorlet{cursfosc}{curs!72!black}%
  \def\fm@logo{\fm@logodir/vertical-#1.pdf}%
  \def\fm@marca{\fm@logodir/marca-#1.pdf}}

% ---------- capçalera de la primera pàgina ----------
% #1 matèria (petit), #2 curs, #3 tipus de document i unitat, #4 títol
\newcommand{\capcalera}[4]{%
  \gdef\fm@cursnom{#2}%
  \par\noindent
  \begin{tcolorbox}[colback=curs,colframe=curs,boxrule=0pt,arc=2.2mm,width=3.5cm,height=2.15cm,
      valign=center,left=3mm,right=2mm,top=1mm,bottom=1mm,nobeforeafter,box align=center]
    \color{curstext}{\footnotesize #1}\par\vspace{1pt}{\fontsize{19}{21}\selectfont\bfseries #2\par}%
  \end{tcolorbox}\hspace{5mm}%
  \begin{minipage}[c]{\dimexpr\linewidth-3.5cm-5mm-2.9cm\relax}
    {\color{cursfosc}\small\bfseries\MakeUppercase{#3}}\par\vspace{2pt}
    {\color{navy}\fontsize{20}{23}\selectfont\bfseries #4\par}
  \end{minipage}\hfill
  \begin{minipage}[c]{2.6cm}\raggedleft\includegraphics[width=2.6cm]{\fm@logo}\end{minipage}\par
  \vspace{2pt}}

% ---------- barres de secció del color del curs ----------
\renewcommand{\seccio}[2]{\par\needspace{\fmneed\baselineskip}\bigskip
  \begin{tcolorbox}[colback=curs,colframe=curs,coltext=curstext,arc=2mm,boxrule=0pt,left=3mm,right=3mm,top=1.2mm,bottom=1.2mm]
  {\large\bfseries #1\quad #2}
  \end{tcolorbox}\vspace{-2pt}}

% ---------- estrelles de dificultat del color del curs (a Batxillerat no se'n posen) ----------
\@ifundefined{estrelles}{}{%
  \renewcommand{\estrelles}[1]{{\color{curs}\ifcase#1\or\faStar\faStar[regular]\faStar[regular]\or\faStar\faStar\faStar[regular]\or\faStar\faStar\faStar\fi}}}

% ---------- capçalera i peu de les pàgines ----------
% primera pàgina: logos institucionals al peu
\fancypagestyle{primera}{\fancyhf{}\renewcommand{\headrulewidth}{0pt}
  \fancyfoot[L]{\raisebox{-0.32\height}{\includegraphics[height=0.78cm]{logo-gva.png}}\hspace{5mm}%
    \raisebox{-0.32\height}{\includegraphics[height=0.78cm]{logo-vermellar.png}}}
  \fancyfoot[R]{\scriptsize\color{muted}\href{https://creativecommons.org/licenses/by-nc-sa/4.0/deed.ca}{\ccbyncsa}\ \ CC BY-NC-SA 4.0 · fentmates.org\qquad\small\thepage}}
% la resta: el símbol a l'esquerra i el curs a la dreta, amb el títol al centre com abans
\fancyhead[L]{\raisebox{-0.3\height}{\includegraphics[height=0.62cm]{\fm@marca}}}
\fancyhead[R]{\small\bfseries\color{cursfosc}\fm@cursnom}
\renewcommand{\headrule}{\color{curs!45}\hrule height \headrulewidth}
\makeatother

%   \fmcurs{1eso}
% i, al començament del document:
%   \capcalera{Matemàtiques}{1r ESO}{Fitxa d'exercicis · Unitat 2}{Divisibilitat}
% Opcional abans de \input: \def\fmneed{14} (espai mínim abans de cada secció).
% =====================================================================
\makeatletter
\def\fm@logodir{../../logo/pdf}
\providecommand{\fmneed}{10}
\def\fm@cursnom{}

% clau del curs, color i color del text damunt del color
\newcommand{\fm@curs}[3]{\@namedef{fm@color@#1}{#2}\@namedef{fm@text@#1}{#3}}
\fm@curs{1eso}{E6B84A}{navy}
\fm@curs{2eso}{3B82F6}{white}
\fm@curs{3eso}{FF7A50}{white}
\fm@curs{4eso-a}{A78BFA}{white}
\fm@curs{4eso-b}{6D28D9}{white}
\fm@curs{1batx-mi}{9B1C1C}{white}
\fm@curs{1batx-ccss}{EF6B6B}{white}
\fm@curs{2batx-mii}{1E6B47}{white}
\fm@curs{2batx-ccss}{2BB3A3}{white}

\newcommand{\fmcurs}[1]{%
  \@ifundefined{fm@color@#1}{\PackageError{fentmates}{Curs desconegut: #1}{}}{}%
  \edef\fm@tmp{\noexpand\definecolor{curs}{HTML}{\@nameuse{fm@color@#1}}%
               \noexpand\colorlet{curstext}{\@nameuse{fm@text@#1}}}\fm@tmp
  \colorlet{cursfosc}{curs!72!black}%
  \def\fm@logo{\fm@logodir/vertical-#1.pdf}%
  \def\fm@marca{\fm@logodir/marca-#1.pdf}}

% ---------- capçalera de la primera pàgina ----------
% #1 matèria (petit), #2 curs, #3 tipus de document i unitat, #4 títol
\newcommand{\capcalera}[4]{%
  \gdef\fm@cursnom{#2}%
  \par\noindent
  \begin{tcolorbox}[colback=curs,colframe=curs,boxrule=0pt,arc=2.2mm,width=3.5cm,height=2.15cm,
      valign=center,left=3mm,right=2mm,top=1mm,bottom=1mm,nobeforeafter,box align=center]
    \color{curstext}{\footnotesize #1}\par\vspace{1pt}{\fontsize{19}{21}\selectfont\bfseries #2\par}%
  \end{tcolorbox}\hspace{5mm}%
  \begin{minipage}[c]{\dimexpr\linewidth-3.5cm-5mm-2.9cm\relax}
    {\color{cursfosc}\small\bfseries\MakeUppercase{#3}}\par\vspace{2pt}
    {\color{navy}\fontsize{20}{23}\selectfont\bfseries #4\par}
  \end{minipage}\hfill
  \begin{minipage}[c]{2.6cm}\raggedleft\includegraphics[width=2.6cm]{\fm@logo}\end{minipage}\par
  \vspace{2pt}}

% ---------- barres de secció del color del curs ----------
\renewcommand{\seccio}[2]{\par\needspace{\fmneed\baselineskip}\bigskip
  \begin{tcolorbox}[colback=curs,colframe=curs,coltext=curstext,arc=2mm,boxrule=0pt,left=3mm,right=3mm,top=1.2mm,bottom=1.2mm]
  {\large\bfseries #1\quad #2}
  \end{tcolorbox}\vspace{-2pt}}

% ---------- estrelles de dificultat del color del curs (a Batxillerat no se'n posen) ----------
\@ifundefined{estrelles}{}{%
  \renewcommand{\estrelles}[1]{{\color{curs}\ifcase#1\or\faStar\faStar[regular]\faStar[regular]\or\faStar\faStar\faStar[regular]\or\faStar\faStar\faStar\fi}}}

% ---------- capçalera i peu de les pàgines ----------
% primera pàgina: logos institucionals al peu
\fancypagestyle{primera}{\fancyhf{}\renewcommand{\headrulewidth}{0pt}
  \fancyfoot[L]{\raisebox{-0.32\height}{\includegraphics[height=0.78cm]{logo-gva.png}}\hspace{5mm}%
    \raisebox{-0.32\height}{\includegraphics[height=0.78cm]{logo-vermellar.png}}}
  \fancyfoot[R]{\scriptsize\color{muted}\href{https://creativecommons.org/licenses/by-nc-sa/4.0/deed.ca}{\ccbyncsa}\ \ CC BY-NC-SA 4.0 · fentmates.org\qquad\small\thepage}}
% la resta: el símbol a l'esquerra i el curs a la dreta, amb el títol al centre com abans
\fancyhead[L]{\raisebox{-0.3\height}{\includegraphics[height=0.62cm]{\fm@marca}}}
\fancyhead[R]{\small\bfseries\color{cursfosc}\fm@cursnom}
\renewcommand{\headrule}{\color{curs!45}\hrule height \headrulewidth}
\makeatother

%   \fmcurs{1eso}
% i, al començament del document:
%   \capcalera{Matemàtiques}{1r ESO}{Fitxa d'exercicis · Unitat 2}{Divisibilitat}
% Opcional abans de \input: \def\fmneed{14} (espai mínim abans de cada secció).
% =====================================================================
\makeatletter
\def\fm@logodir{../../logo/pdf}
\providecommand{\fmneed}{10}
\def\fm@cursnom{}

% clau del curs, color i color del text damunt del color
\newcommand{\fm@curs}[3]{\@namedef{fm@color@#1}{#2}\@namedef{fm@text@#1}{#3}}
\fm@curs{1eso}{E6B84A}{navy}
\fm@curs{2eso}{3B82F6}{white}
\fm@curs{3eso}{FF7A50}{white}
\fm@curs{4eso-a}{A78BFA}{white}
\fm@curs{4eso-b}{6D28D9}{white}
\fm@curs{1batx-mi}{9B1C1C}{white}
\fm@curs{1batx-ccss}{EF6B6B}{white}
\fm@curs{2batx-mii}{1E6B47}{white}
\fm@curs{2batx-ccss}{2BB3A3}{white}

\newcommand{\fmcurs}[1]{%
  \@ifundefined{fm@color@#1}{\PackageError{fentmates}{Curs desconegut: #1}{}}{}%
  \edef\fm@tmp{\noexpand\definecolor{curs}{HTML}{\@nameuse{fm@color@#1}}%
               \noexpand\colorlet{curstext}{\@nameuse{fm@text@#1}}}\fm@tmp
  \colorlet{cursfosc}{curs!72!black}%
  \def\fm@logo{\fm@logodir/vertical-#1.pdf}%
  \def\fm@marca{\fm@logodir/marca-#1.pdf}}

% ---------- capçalera de la primera pàgina ----------
% #1 matèria (petit), #2 curs, #3 tipus de document i unitat, #4 títol
\newcommand{\capcalera}[4]{%
  \gdef\fm@cursnom{#2}%
  \par\noindent
  \begin{tcolorbox}[colback=curs,colframe=curs,boxrule=0pt,arc=2.2mm,width=3.5cm,height=2.15cm,
      valign=center,left=3mm,right=2mm,top=1mm,bottom=1mm,nobeforeafter,box align=center]
    \color{curstext}{\footnotesize #1}\par\vspace{1pt}{\fontsize{19}{21}\selectfont\bfseries #2\par}%
  \end{tcolorbox}\hspace{5mm}%
  \begin{minipage}[c]{\dimexpr\linewidth-3.5cm-5mm-2.9cm\relax}
    {\color{cursfosc}\small\bfseries\MakeUppercase{#3}}\par\vspace{2pt}
    {\color{navy}\fontsize{20}{23}\selectfont\bfseries #4\par}
  \end{minipage}\hfill
  \begin{minipage}[c]{2.6cm}\raggedleft\includegraphics[width=2.6cm]{\fm@logo}\end{minipage}\par
  \vspace{2pt}}

% ---------- barres de secció del color del curs ----------
\renewcommand{\seccio}[2]{\par\needspace{\fmneed\baselineskip}\bigskip
  \begin{tcolorbox}[colback=curs,colframe=curs,coltext=curstext,arc=2mm,boxrule=0pt,left=3mm,right=3mm,top=1.2mm,bottom=1.2mm]
  {\large\bfseries #1\quad #2}
  \end{tcolorbox}\vspace{-2pt}}

% ---------- estrelles de dificultat del color del curs (a Batxillerat no se'n posen) ----------
\@ifundefined{estrelles}{}{%
  \renewcommand{\estrelles}[1]{{\color{curs}\ifcase#1\or\faStar\faStar[regular]\faStar[regular]\or\faStar\faStar\faStar[regular]\or\faStar\faStar\faStar\fi}}}

% ---------- capçalera i peu de les pàgines ----------
% primera pàgina: logos institucionals al peu
\fancypagestyle{primera}{\fancyhf{}\renewcommand{\headrulewidth}{0pt}
  \fancyfoot[L]{\raisebox{-0.32\height}{\includegraphics[height=0.78cm]{logo-gva.png}}\hspace{5mm}%
    \raisebox{-0.32\height}{\includegraphics[height=0.78cm]{logo-vermellar.png}}}
  \fancyfoot[R]{\scriptsize\color{muted}\href{https://creativecommons.org/licenses/by-nc-sa/4.0/deed.ca}{\ccbyncsa}\ \ CC BY-NC-SA 4.0 · fentmates.org\qquad\small\thepage}}
% la resta: el símbol a l'esquerra i el curs a la dreta, amb el títol al centre com abans
\fancyhead[L]{\raisebox{-0.3\height}{\includegraphics[height=0.62cm]{\fm@marca}}}
\fancyhead[R]{\small\bfseries\color{cursfosc}\fm@cursnom}
\renewcommand{\headrule}{\color{curs!45}\hrule height \headrulewidth}
\makeatother

\fmcurs{2eso}

\begin{document}
\thispagestyle{primera}

\capcalera{Matemàtiques}{2n ESO}{Fitxa d'exercicis · Unitat 3}{Introducció a l'àlgebra.\\ Polinomis}
\vspace{10pt}
Nom i cognoms: \rule{8.4cm}{0.4pt}\hfill Grup: \rule{1.6cm}{0.4pt}\hfill Data: \rule{2.2cm}{0.4pt}
\vspace{8pt}

\begin{tcolorbox}[colback=soft,colframe=soft,arc=2mm,left=3mm,right=3mm,top=1.5mm,bottom=1.5mm]
\small Exercicis graduats per nivells de dificultat:\quad \estrelles{1}\ bàsic \qquad \estrelles{2}\ mitjà \qquad \estrelles{3}\ avançat.\\[2pt] Resol-los al quadern mostrant tots els passos i sense calculadora. Dona els resultats reduïts i ordenats. Les seccions segueixen l'ordre dels apunts de la unitat (3.1--3.8).
\end{tcolorbox}

% =====================================================
\seccio{1}{L'àlgebra: el llenguatge algebraic}

\ex{1} Escriu en llenguatge algebraic (fes servir $x$ per al nombre):\begin{tasks}[after-item-skip=3pt](2)
  \task El triple d'un nombre.
  \task Un nombre menys 8.
  \task La meitat d'un nombre.
  \task El quadrat d'un nombre més 1.
  \task El doble d'un nombre menys 5.
  \task El doble de la diferència d'un nombre i 5.
  \task El següent d'un nombre.
  \task La suma d'un nombre i el seu quadrat.
  \task El cub d'un nombre.
  \task La quarta part d'un nombre menys 3.
\end{tasks}\recorda{l'ordre de les paraules importa: «el doble d'un nombre menys 5» és $2x-5$, però «el doble de la diferència d'un nombre i 5» és $2(x-5)$.}

\ex{1} Escriu amb paraules què expressa cada expressió algebraica:\begin{tasks}(4)
  \task $x+5$
  \task $4x$
  \task $x^2-3$
  \task $\dfrac{x}{3}$
  \task $2x+1$
  \task $(x+2)^2$
  \task $x-y$
  \task $10-x$
\end{tasks}

\ex{1} Escriu l'expressió algebraica de: a) el perímetre d'un quadrat de costat $x$; b) l'àrea d'un rectangle de base $x$ i altura $5$; c) l'edat d'una persona d'ací a 7 anys si ara en té $x$; d) el preu de $x$ quilos de pomes a $1{,}80$ € el quilo.

\ex{1} Calcula el valor numèric de cada expressió per al valor de $x$ indicat:\begin{tasks}[after-item-skip=5pt](3)
  \task $3x+5$ \ per a $x=2$
  \task $7-2x$ \ per a $x=-3$
  \task $x^{2}-4x$ \ per a $x=5$
  \task $2x^{2}+x-1$ \ per a $x=-2$
  \task $-x^{2}+6$ \ per a $x=-3$
  \task $x^{3}-x$ \ per a $x=-2$
\end{tasks}\recorda{els nombres negatius se substitueixen entre parèntesis: si $x=-3$, \ $x^2=(-3)^2=9$ però $-x^2=-(-3)^2=-9$.}

\ex{2} Calcula el valor numèric:\begin{tasks}[after-item-skip=5pt](2)
  \task $2a+3b$ \ per a $a=4,\ b=-1$
  \task $a^{2}-b^{2}$ \ per a $a=-3,\ b=2$
  \task $5xy-x$ \ per a $x=2,\ y=-1$
  \task $x^{2}y+xy^{2}$ \ per a $x=-1,\ y=3$
\end{tasks}

\ex{2} Una companyia de telefonia cobra una quota de 9 € al mes més $0{,}05$ € per cada minut de trucada. Escriu l'expressió del que paguem en un mes si parlem $x$ minuts i calcula quant paguem si parlem 120 minuts. I si parlem 3 hores?

\ex{2} Amb escuradents formem una filera de triangles: la figura 1 té 3 escuradents, la 2 en té 5, la 3 en té 7… Quants en té la figura 10? Escriu una expressió per als escuradents de la figura $n$. Quina figura té 51 escuradents?

\ex{3} Escriu tres nombres consecutius qualssevol fent servir $x$ i demostra que la seua suma sempre és múltiple de 3. És cert també per a la suma de quatre nombres consecutius?


% =====================================================
\seccio{2}{Monomis: parts i valor numèric}

\ex{1} Completa la taula:\par\smallskip{\centering\renewcommand{\arraystretch}{1.35}\begin{tabular}{|c|C{2.6cm}|C{2.6cm}|C{1.6cm}|}\hline Monomi & Coeficient & Part literal & Grau\\\hline $-7x^3$ & & & \\\hline $xy^2$ & & & \\\hline $\tfrac{2}{3}a^2b$ & & & \\\hline $-x$ & & & \\\hline $5x^2y^3z$ & & & \\\hline $12$ & & & \\\hline\end{tabular}\par}\recorda{el grau d'un monomi és la \textbf{suma} dels exponents de totes les lletres; si no veus el coeficient, és 1 (o $-1$).}

\ex{1} Quines d'aquestes expressions són monomis? Explica per què les altres no ho són.\begin{tasks}(4)
  \task $3x^4$
  \task $\dfrac{5}{x}$
  \task $x+y$
  \task $-\dfrac{x^2}{4}$
  \task $\sqrt{x}$
  \task $-2ab^3$
  \task $x^{-1}$
  \task $9$
\end{tasks}

\ex{1} Agrupa els monomis semblants: \ $4x^2,\ \ -3xy,\ \ 2x,\ \ -x^2,\ \ \tfrac{1}{2}xy,\ \ 7x,\ \ 5x^2y,\ \ -x,\ \ \tfrac{3}{4}x^2,\ \ x^2y$.

\ex{2} Calcula el valor numèric de cada monomi:\begin{tasks}[after-item-skip=5pt](2)
  \task $-3x^{2}$ \ per a $x=4$
  \task $2x^{3}$ \ per a $x=-2$
  \task $-a^{2}b$ \ per a $a=-3,\ b=2$
  \task $4xy^{2}$ \ per a $x=\tfrac{1}{2},\ y=-3$
\end{tasks}

\ex{2} Escriu un monomi:\begin{tasks}[after-item-skip=3pt](1)
  \task semblant a $-6x^2y$ amb coeficient positiu;
  \task de grau 5 amb dues lletres i coeficient $-\tfrac{1}{2}$;
  \task oposat a $3ab^2$;
  \task de grau 0.
\end{tasks}

\ex{3} Troba $m$ i $n$ perquè els monomis $5x^{m}y^{3}$ i $-2x^{4}y^{n+1}$ siguen semblants. Quin grau tenen?

\ex{3} Escriu un monomi de grau 4, en les lletres $x$ i $y$, amb coeficient 2, que valga 54 per a $x=3$, $y=1$.


% =====================================================
\seccio{3}{Operacions amb monomis}

\ex{1} Redueix:\begin{tasks}[after-item-skip=6pt](3)
  \task $5x+3x-2x$
  \task $4x^{2}-9x^{2}$
  \task $-a+6a-8a$
  \task $3xy-xy+5xy$
  \task $x^{3}+x^{3}+x^{3}$
  \task $7ab^{2}-7ab^{2}$
\end{tasks}\recorda{només se sumen i resten monomis \textbf{semblants}: sumem els coeficients i \textbf{deixem la mateixa part literal} ($5x^2+3x^2=8x^2$, mai $8x^4$).}

\ex{1} Redueix quan siga possible. Si no es pot, explica per què:\begin{tasks}(4)
  \task $3x^{2}+5x$
  \task $4x-x^{2}+2x$
  \task $6ab-2a+ab$
  \task $2x^{3}-x^{2}+4x^{3}-x$
\end{tasks}

\ex{1} Multiplica:\begin{tasks}[after-item-skip=6pt](3)
  \task $3x^{2}\cdot 5x^{4}$
  \task $(-2a)\cdot (7a^{3})$
  \task $(-4xy)\cdot (-3x^{2}y)$
  \task $6x^{2}\cdot (\tfrac{1}{3})x$
  \task $x^{3}y\cdot (-xy^{4})$
  \task $2a\cdot 3b\cdot (-a^{2})$
\end{tasks}\recorda{multipliquem els coeficients i \textbf{sumem} els exponents de cada lletra: $x^m\cdot x^n=x^{m+n}$.}

\ex{1} Divideix:\begin{tasks}[after-item-skip=6pt](3)
  \task $12x^{5}:(4x^{2})$
  \task $(-20a^{4}):(5a)$
  \task $18x^{3}y^{2}:(-6xy)$
  \task $(-9x^{6}):(-9x^{6})$
  \task $7x^{4}y:(14x^{2})$
  \task $(-15a^{3}b^{4}):(3ab^{4})$
\end{tasks}\recorda{dividim els coeficients i \textbf{restem} els exponents: $x^m:x^n=x^{m-n}$; i $x^0=1$.}

\ex{2} Calcula les potències:\begin{tasks}[after-item-skip=6pt](3)
  \task $(2x^{3})^{2}$
  \task $(-3a^{2})^{3}$
  \task $(x^{2}y^{3})^{4}$
  \task $(-2xy^{2})^{4}$
  \task $(-a^{3}b)^{5}$
  \task $(\tfrac{1}{2}x^{2})^{3}$
\end{tasks}\recorda{elevem el coeficient i \textbf{multipliquem} els exponents: $(-2x^3y)^2=(-2)^2x^{6}y^{2}=4x^6y^2$. Compte amb el signe si l'exponent és senar.}

\ex{2} Redueix agrupant els termes semblants:\begin{tasks}[after-item-skip=6pt](2)
  \task $5x^{2}-3x+2-x^{2}+7x-6$
  \task $-2a+4b-a-3b+5a$
  \task $x^{3}-2x+x^{2}-4x^{3}+3x+x^{2}$
  \task $3xy-x^{2}+2xy+4x^{2}-y$
\end{tasks}

\ex{2} Quins d'aquests quocients són monomis? Calcula'ls.\begin{tasks}(4)
  \task $8x^{4}:(2x)$
  \task $6x^{2}:(3x^{5})$
  \task $10a^{3}b:(5ab)$
  \task $4xy:(2y^{2})$
\end{tasks}

\ex{2} Opera i redueix (respecta la jerarquia de les operacions):\begin{tasks}[after-item-skip=6pt](2)
  \task $3x\cdot 4x^{2}-5x^{3}$
  \task $2a^{2}\cdot (-3a)+10a^{4}:(2a)$
  \task $(2x)^{3}-x\cdot 6x^{2}$
  \task $15x^{4}:(3x^{2})-(2x)^{2}$
  \task $-x^{2}\cdot x^{3}+(x^{2})^{2}\cdot 3x$
  \task $(-3x)^{2}\cdot x-4x^{3}$
\end{tasks}

\ex{3} Opera i redueix:\begin{tasks}[after-item-skip=6pt](2)
  \task $(-2x^{2}y)^{3}:(4x^{4}y^{2})$
  \task $(3a^{2}b)^{2}\cdot (-ab)^{3}:(9a^{5}b^{4})$
  \task $(-x^{2})^{3}+2x^{4}\cdot (-x)^{2}-(2x^{3})^{2}:4$
  \task $(2x)^{4}:(-x)^{2}-3x\cdot (-2x)$
\end{tasks}

\ex{3} Un ortoedre té les arestes $2x$, $3x$ i $x^2$ cm. Escriu com a monomis o polinomis: a) el volum; b) l'àrea total. c) Què li passa al volum si totes les arestes es fan el doble de llargues?


% =====================================================
\seccio{4}{Polinomis: parts}

\ex{1} Completa la taula:\par\smallskip{\centering\renewcommand{\arraystretch}{1.4}\small\begin{tabular}{|c|C{3.8cm}|C{1cm}|C{1.7cm}|C{1.7cm}|C{1.5cm}|}\hline Polinomi & Ordenat & Grau & Coef. principal & Terme indep. & Termes\\\hline $3x^{2}-5x+1$ & & & & & \\\hline $7-2x^{4}+x$ & & & & & \\\hline $x^{3}-8$ & & & & & \\\hline $-x^{5}+4x^{3}-x^{2}+2x$ & & & & & \\\hline $6x$ & & & & & \\\hline $4x^{2}-x+3x^{3}-2$ & & & & & \\\hline\end{tabular}\par}\recorda{el \textbf{grau} d'un polinomi és el grau més alt dels seus termes; el \textbf{coeficient principal}, el del terme de grau més alt; el \textbf{terme independent}, el que no té lletra (amb el seu signe).}

\ex{1} Redueix i ordena aquests polinomis. Quin grau té cada un?\begin{tasks}[after-item-skip=6pt](2)
  \task $2x^{2}+3x-x^{2}+5-x$
  \task $x^{3}-4x+2x^{3}+x^{2}-3x^{3}+4x$
  \task $5-x+3x^{2}-7+x-x^{2}$
  \task $-x^{4}+2x+x^{4}-3x^{2}+x^{3}$
\end{tasks}

\ex{1} Quins d'aquests polinomis són complets? Escriu els que no ho són posant zeros als termes que falten.\begin{tasks}(2)
  \task $x^{3}+2x^{2}-x+5$
  \task $4x^{2}-1$
  \task $x^{4}-x^{3}+x^{2}-x+1$
  \task $2x^{3}+x$
\end{tasks}

\ex{2} Calcula el valor numèric de $P(x)=x^2-3x+2$ i de $Q(x)=2x^3-x^2+4$ per als valors indicats:\begin{tasks}[after-item-skip=5pt](4)
  \task $P(0)$
  \task $P(1)$
  \task $P(-1)$
  \task $P(2)$
  \task $Q(0)$
  \task $Q(1)$
  \task $Q(-2)$
  \task $Q\!\left(\tfrac{1}{2}\right)$
\end{tasks}

\ex{2} Escriu un polinomi que complisca cada condició: a) de grau 3 i amb tres termes; b) de grau 2, complet i amb terme independent $-4$; c) de grau 4, amb coeficient principal $-2$ i sense terme independent.

\ex{3} a) Troba $k$ perquè $P(x)=x^3-2x^2+kx+6$ complisca $P(2)=0$. \ b) Troba $a$ si $Q(x)=ax^2+3x-1$ i $Q(2)=13$.

\ex{3} Un polinomi de grau 2 té el terme independent igual a 5, el coeficient de $x$ igual a $-4$ i compleix $P(1)=3$. Quin és? Quant val $P(-1)$?


% =====================================================
\seccio{5}{Suma i resta de polinomis}

\ex{1} Donats $P(x)=3x^2-5x+2$, \ \ $Q(x)=x^3-2x^2+4x-1$ \ \ i \ \ $R(x)=-x^2+x-6$, calcula:\begin{tasks}(4)
  \task $P+Q$
  \task $P+R$
  \task $Q+R$
  \task $P+Q+R$
\end{tasks}

\ex{1} Amb els mateixos polinomis, calcula:\begin{tasks}(4)
  \task $P-Q$
  \task $Q-P$
  \task $P-R$
  \task $R-Q$
\end{tasks}\recorda{per a restar, canviem el signe de \textbf{tots} els termes del polinomi que restem: $-(x^3-2x^2+4x-1)=-x^3+2x^2-4x+1$.}

\ex{1} Escriu l'oposat de cada polinomi i comprova que la seua suma és 0:\begin{tasks}(3)
  \task $2x^{2}-7x+1$
  \task $-x^{3}+x-5$
  \task $4x^{4}-3x^{2}$
\end{tasks}

\ex{2} Amb els polinomis de l'exercici 33, calcula:\begin{tasks}(4)
  \task $P+Q-R$
  \task $P-(Q+R)$
  \task $Q-P-R$
  \task $2P-R$
\end{tasks}

\ex{2} Opera i redueix:\begin{tasks}[after-item-skip=6pt](1)
  \task $(x^{2}-4x+3)-(2x^{2}+x-5)+(x^{2}-1)$
  \task $(5x^{3}-2x)-(x^{3}+x^{2}-2x+4)$
  \task $-(x^{2}-3)+(4x^{2}-x)-(2x+7)$
  \task $(2x^{2}y-xy+3)-(x^{2}y+4xy-1)$
\end{tasks}

\ex{3} Troba el polinomi $X(x)$ en cada cas:\begin{tasks}[after-item-skip=4pt](1)
  \task $X(x)+(2x^2-3x+1)=5x^2+x-4$
  \task $(x^3-x+2)-X(x)=3x^2+x$
\end{tasks}

\ex{3} Els costats d'un pentàgon fan $x^2+1$, $2x$, $x^2-x+3$, $3x+2$ i $2x^2-4$ cm. Escriu-ne el perímetre i calcula'l per a $x=4$. Quant més llarg és el costat més llarg que el més curt quan $x=4$?


% =====================================================
\seccio{6}{Producte de polinomis}

\ex{1} Multiplica:\begin{tasks}[after-item-skip=6pt](3)
  \task $3(x^{2}-2x+5)$
  \task $-2(4x^{3}-x+1)$
  \task $5(x^{2}-3x)$
  \task $-(2x^{2}-x-7)$
  \task $\tfrac{1}{2}(4x^{2}-6x+2)$
  \task $-4(-x^{3}+2x^{2}-3)$
\end{tasks}\recorda{propietat distributiva: multipliquem el nombre o el monomi per \textbf{cada} terme del polinomi.}

\ex{1} Multiplica:\begin{tasks}[after-item-skip=6pt](3)
  \task $2x(x^{2}-3x+4)$
  \task $-3x^{2}(2x-5)$
  \task $x^{3}(x^{2}+x-1)$
  \task $-4x(x^{3}-2x+3)$
  \task $5xy(x-2y+1)$
  \task $-x^{2}(-x^{2}+3x-6)$
\end{tasks}

\ex{1} Multiplica i redueix:\begin{tasks}[after-item-skip=6pt](3)
  \task $(x+2)(x+5)$
  \task $(x-3)(x+4)$
  \task $(x-1)(x-6)$
  \task $(2x+1)(x-3)$
  \task $(3x-2)(2x+5)$
  \task $(x^{2}+1)(x-2)$
\end{tasks}\recorda{cada terme del primer per cada terme del segon: $(x+3)(2x-5)=2x^2-5x+6x-15=2x^2+x-15$.}

\ex{2} Multiplica i redueix:\begin{tasks}[after-item-skip=6pt](2)
  \task $(x+2)(x^{2}-3x+1)$
  \task $(2x-1)(x^{2}+4x-3)$
  \task $(x^{2}-x+2)(3x+1)$
  \task $(x^{2}+2x-1)(x^{2}-2)$
\end{tasks}

\ex{2} Opera i redueix:\begin{tasks}[after-item-skip=6pt](2)
  \task $3(x+2)-2(x-1)$
  \task $x(x+4)-2(x^{2}-3x)$
  \task $(x+1)(x-2)+(x+3)(x-1)$
  \task $2x(x-5)-(x+2)(x-3)$
\end{tasks}

\ex{2} Escriu, com a polinomi reduït, l'àrea de: a) un rectangle de costats $x+4$ i $2x-1$; b) un quadrat de costat $x+3$ al qual llevem un quadrat de costat $x$; c) un triangle de base $2x+6$ i altura $x$.

\ex{3} Multiplica (fes primer el producte de dos factors i després multiplica pel tercer):\begin{tasks}[after-item-skip=6pt](2)
  \task $(x+1)(x-1)(x+3)$
  \task $(x-2)(2x+1)(x+2)$
\end{tasks}

\ex{3} Opera i redueix:\begin{tasks}[after-item-skip=6pt](1)
  \task $(2x-1)(x+3)-(x-2)(x^{2}+x)+x^{3}$
  \task $(x^{2}-2x+3)(x+1)-x(x^{2}-x+1)$
\end{tasks}


% =====================================================
\seccio{7}{Identitats notables}

\ex{1} Desenvolupa (quadrat d'una suma):\begin{tasks}[after-item-skip=6pt](4)
  \task $(x+3)^{2}$
  \task $(x+7)^{2}$
  \task $(2x+1)^{2}$
  \task $(a+5b)^{2}$
\end{tasks}\recorda{$(a+b)^2=a^2+2ab+b^2$: \ el quadrat del primer, \textbf{més el doble del primer pel segon}, més el quadrat del segon.}

\ex{1} Desenvolupa (quadrat d'una diferència):\begin{tasks}[after-item-skip=6pt](4)
  \task $(x-4)^{2}$
  \task $(x-1)^{2}$
  \task $(3x-2)^{2}$
  \task $(5-x)^{2}$
\end{tasks}\recorda{$(a-b)^2=a^2-2ab+b^2$.}

\ex{1} Desenvolupa (suma per diferència):\begin{tasks}[after-item-skip=6pt](4)
  \task $(x+5)(x-5)$
  \task $(x-9)(x+9)$
  \task $(2x+3)(2x-3)$
  \task $(4+x)(4-x)$
\end{tasks}\recorda{$(a+b)(a-b)=a^2-b^2$: suma per diferència és igual a diferència de quadrats.}

\ex{2} Desenvolupa fent servir les identitats notables:\begin{tasks}[after-item-skip=6pt](3)
  \task $(3x+4)^{2}$
  \task $(2x-5y)^{2}$
  \task $(x^{2}+3)^{2}$
  \task $(x^{2}-2x)^{2}$
  \task $(3a+2b)(3a-2b)$
  \task $(\tfrac{1}{2}x-2)^{2}$
\end{tasks}

\ex{2} Escriu com a quadrat d'un binomi o com a producte d'una suma per una diferència:\begin{tasks}[after-item-skip=6pt](3)
  \task $x^{2}+10x+25$
  \task $x^{2}-6x+9$
  \task $x^{2}-16$
  \task $4x^{2}+4x+1$
  \task $9x^{2}-25$
  \task $x^{2}-2x+1$
\end{tasks}

\ex{2} Completa les caselles:\begin{tasks}[after-item-skip=6pt](2)
  \task $(x+\square)^2=x^2+\square+16$
  \task $(x-\square)^2=x^2-10x+\square$
  \task $(\square+3)^2=4x^2+\square+9$
  \task $(x+\square)(x-\square)=x^2-49$
\end{tasks}

\ex{3} Simplifica:\begin{tasks}[after-item-skip=6pt](1)
  \task $(x+2)^{2}-(x-2)^{2}$
  \task $(2x-1)^{2}-(2x+1)(2x-1)$
  \task $(x+3)^{2}-2(x-1)(x+1)-x(x+6)$
\end{tasks}

\ex{3} Calcula sense calculadora, fent servir les identitats notables:\begin{tasks}(4)
  \task $41^2$
  \task $29^2$
  \task $103\cdot97$
  \task $55^2-45^2$
\end{tasks}

\ex{3} a) Demostra que $(a+b)^2-(a-b)^2=4ab$. \ b) Fes-la servir per a calcular $(x+5)^2-(x-5)^2$ i $\ 1001^2-999^2$.


% =====================================================
\seccio{8}{Extracció de factor comú}

\ex{1} Trau factor comú:\begin{tasks}[after-item-skip=6pt](3)
  \task $3x+12$
  \task $5x^{2}-10x$
  \task $x^{3}+x^{2}$
  \task $8x-12$
  \task $7x^{2}+14x^{3}$
  \task $6x^{3}-4x^{2}+2x$
\end{tasks}\recorda{el factor comú és el MCD dels coeficients per les lletres comunes a tots els termes amb el \textbf{menor} exponent. Dins del parèntesi va el que queda en dividir cada terme entre el factor comú (si un terme és igual al factor comú, hi queda un 1).}

\ex{1} Calcula de la manera més ràpida, traient factor comú:\begin{tasks}(4)
  \task $34\cdot 18+34\cdot 82$
  \task $125\cdot 7-125\cdot 3$
  \task $9\cdot 47+9\cdot 53$
  \task $56^2-56\cdot 46$
\end{tasks}

\ex{2} Trau factor comú:\begin{tasks}[after-item-skip=6pt](2)
  \task $4a^{2}b-6ab^{2}$
  \task $9x^{4}-6x^{3}+3x^{2}$
  \task $10x^{3}y^{2}+15x^{2}y^{3}$
  \task $x^{5}-x^{4}+x^{3}-x^{2}$
  \task $-8x^{3}-12x$
  \task $12a^{2}b^{3}-18a^{3}b^{2}+6a^{2}b^{2}$
\end{tasks}

\ex{2} Completa les caselles:\begin{tasks}[after-item-skip=6pt](2)
  \task $6x^3-9x^2=3x^2(\square-\square)$
  \task $10x^4+5x^2=5x^2(\square+\square)$
  \task $x^2y-xy^2=\square(x-y)$
  \task $14a^3-21a^2b=\square(2a-3b)$
\end{tasks}

\ex{3} Factoritza tant com pugues (primer el factor comú i després una identitat notable):\begin{tasks}[after-item-skip=6pt](2)
  \task $2x^{3}-8x$
  \task $3x^{3}+6x^{2}+3x$
  \task $x^{4}-9x^{2}$
  \task $5x^{3}-20x^{2}+20x$
\end{tasks}

\ex{3} Trau factor comú (el factor comú és un binomi):\begin{tasks}[after-item-skip=6pt](3)
  \task $x(x+2)+3(x+2)$
  \task $2x(x-5)-(x-5)$
  \task $(x+1)^{2}+4(x+1)$
\end{tasks}

\ex{3} Sense calculadora:\begin{tasks}[after-item-skip=4pt](1)
  \task $99^2+99$
  \task $13\cdot87+13^2$
  \task Demostra que $n^2+n$ sempre és parell per a qualsevol nombre natural $n$.
\end{tasks}

\vfill
\begin{center}\color{muted}\small\faCheckCircle[regular]\ \ Ànim! Comprova els teus resultats amb les eines interactives dels apunts de la unitat.\end{center}

\end{document}
